\clearpage\pagestyle{otMain}\setlength{\otBlocEspacement}{2.5mm}\fontsize{11}{15}\selectfont\corrigesChapitre{2}{Additionner, soustraire, multiplier}
\begin{blocLivre}{cor-2-1}\begin{corrige}{1}{Deux sommes posées}\demarche{Aligner les rangs et estimer.}
$38\,675+4\,850=\res{43\,525}$ et $256\,908+87\,314=\res{344\,222}$. Estimations : $43\,500$ et $344\,000$ environ.
\end{corrige}\end{blocLivre}
\begin{blocLivre}{cor-2-2}\begin{corrige}{2}{Deux différences posées}\demarche{Contrôler par l’addition.}
$70\,200-2\,460=\res{67\,740}$ ; contrôle : $67\,740+2\,460=70\,200$. Et $500\,000-123\,456=\res{376\,544}$.
\end{corrige}\end{blocLivre}
\begin{blocLivre}{cor-2-3}\begin{corrige}{3}{Presque cent}\demarche{Passer par la centaine la plus proche.}
$457+100-1=\res{556}$ ; $623-100+2=\res{525}$ ; $350+200-2=\res{548}$.
\end{corrige}\end{blocLivre}
\begin{blocLivre}{cor-2-4}\begin{corrige}{4}{Des nombres ronds}\demarche{Multiplier les chiffres, compter les zéros.}
$3\times4=12$ puis trois zéros : $\res{12\,000}$. $7\times8=56$ puis trois zéros : $\res{56\,000}$. $6\times5=30$ puis trois zéros : $\res{30\,000}$ ; le zéro de 30 s’ajoute.
\end{corrige}\end{blocLivre}
\begin{blocLivre}{cor-2-5}\begin{corrige}{5}{Un produit posé}\demarche{Séparer unités et dizaines du multiplicateur.}
$485\times6=2\,910$ ; $485\times30=14\,550$ ; somme : $\res{17\,460}$. Estimation : $500\times35=17\,500$.
\end{corrige}\end{blocLivre}
\begin{blocLivre}{cor-2-6}\begin{corrige}{6}{Les parenthèses d’abord}\demarche{Traiter l’intérieur des parenthèses en premier.}
$5\times(12-4)=5\times8=\res{40}$ ; $5\times12-4=60-4=\res{56}$. Les parenthèses imposent de soustraire avant de multiplier : l’ordre des actions change le résultat.
\end{corrige}\end{blocLivre}
\begin{blocLivre}{cor-2-7}\begin{corrige}{7}{Deux paires}\demarche{Commencer par la parenthèse la plus intérieure.}
$50\times4=\res{200}$. Pour la seconde : la parenthèse intérieure donne 8, puis $(20-8)\times2=12\times2=\res{24}$.
\end{corrige}\end{blocLivre}
\begin{blocLivre}{cor-2-8}\begin{corrige}{8}{Traduire une phrase}\demarche{Encadrer la quantité décrite avant de multiplier.}
La différence vaut $40-13$ ; son triple s’écrit $\res{3\times(40-13)}=3\times27=\res{81}$. Sans parenthèses, $3\times40-13=107$ dirait autre chose.
\end{corrige}\end{blocLivre}
\begin{blocLivre}{cor-2-9}\begin{corrige}{9}{La bonne expression}\demarche{Chercher ce que paie une seule personne.}
Chacun dépense $4+2$ euros : il faut $\res{3\times(4+2)}=18$ euros. L’expression $3\times4+2=14$ compterait une seule boisson.
\end{corrige}\end{blocLivre}
\begin{blocLivre}{cor-2-10}\begin{corrige}{10}{Compléter les étapes}\demarche{Choisir l’opération qui défait l’autre.}
$56\times\res{100}=5\,600$ ; $12\,800+4\,500=\res{17\,300}$ ; $4\,900:7=\res{700}$.
\end{corrige}\end{blocLivre}
\begin{blocLivre}{cor-2-11}\begin{corrige}{11}{Le festival}\demarche{Additionner avant de comparer à la capacité.}
Billets vendus : $12\,760+9\,845=22\,605$. Places libres : $25\,000-22\,605=\res{2\,395}$. L’estimation $12\,800+9\,800=22\,600$ confirme la somme.
\end{corrige}\end{blocLivre}
\begin{blocLivre}{cor-2-12}\begin{corrige}{12}{La commande de la mairie}\demarche{Calculer chaque produit avant la somme.}
Tables : $485\times36=17\,460$ euros. Chaises : $59\times144=8\,496$ euros. Total : $17\,460+8\,496=\res{25\,956}$ euros.
\end{corrige}\end{blocLivre}
\begin{blocLivre}{cor-2-13}\begin{corrige}{13}{Une facture à vérifier}\demarche{Estimer avant de refaire le calcul exact.}
Estimation : $24\times18$ est proche de $25\times18=450$ : 532 semble trop grand. Calcul exact : $24\times18=\res{432}$ euros ; la facture $\res{\text{est fausse}}$ de 100 euros.
\end{corrige}\end{blocLivre}
\begin{blocLivre}{cor-2-14}\begin{corrige}{14}{Le double piège}\demarche{Refaire chaque calcul en respectant les priorités.}
$2\times15+5=30+5=\res{35}$ : c’est Zoé qui se trompe. $2\times(15+5)=2\times20=\res{40}$ : Tom a raison. Même écriture sans parenthèses, autre sens.
\end{corrige}\end{blocLivre}
\begin{blocLivre}{cor-2-15}\begin{corrige}{15}{Le tour de magie}\demarche{Suivre ce que devient le nombre choisi.}
Avec 40 : $40+25-17-40=\res8$. Avec 358 : $358+25-17-358=\res8$. Le nombre choisi s’ajoute puis se retire : il disparaît du bilan, qui vaut toujours $25-17=8$.
\end{corrige}\end{blocLivre}
